\selectlanguage{italian}
\Language{it}{Italiano}{Guida all'uso}
\firsttopic{Aprire, sospendere, tornare}
Inkline esegue una shell locale su reMarkable 2. Usa Type Folio o una tastiera
USB, lavora sul tablet oppure collegati con SSH o Goblin Mosh. Testo e immagini
sono visualizzati in scala di grigi.

\begin{notice}
\key{Opt+RightAlt+T} apre Inkline direttamente dal tablet.\par
\key{Pulsante di accensione:} premi e rilascia per sospendere; premi di nuovo
per riattivare. Le shell e i buffer dell'editor rimangono in RAM.
\end{notice}
La sospensione mette in pausa il tablet senza chiudere Inkline né salvare le
impostazioni sulla memoria flash. Le connessioni di rete potrebbero doversi
ristabilire. La pressione che riattiva non provoca un'altra sospensione.
Dopo un aggiornamento, chiudi e riapri Inkline per attivare la nuova gestione del pulsante.

Per tornare ai quaderni, tocca \key{Quit} oppure premi \key{Ctrl+Shift+Q} e conferma.
\code{exit} chiude un solo terminale; chiudendo l'ultimo tornano i quaderni.
La barra inferiore contiene Escape, batteria, Quit e Settings.
\code{clear} cancella la breve guida iniziale e il piccolo logo Goblin.

Se un'app avviata da una scorciatoia occupa lo schermo, \key{Opt+RightAlt+T} torna a Inkline.
Per il recupero d'emergenza, tieni premuto \key{Opt+RightAlt+Backspace per due secondi}.
Il recupero chiude tutti i terminali e perde il lavoro non salvato in quelle
sessioni. Dopo il riavvio del tablet, l'interfaccia iniziale resta quella dei quaderni.

RightAlt indica il tasto Alt/Opt a destra della barra spaziatrice. Tienilo premuto insieme al tasto Opt separato, poi premi T, senza Ctrl. Super da solo resta disponibile per le tue scorciatoie.

\ProjectLinks

\topic{Tasti e slot dei terminali}
Type Folio ha un tasto \key{Opt} separato fra Ctrl e Alt e un tasto
\key{Alt/Opt destro}. Svolgono funzioni diverse.

\begin{keytable}
\key{Tasti} & \key{Azione} \\ \midrule
Opt separato + 1--0 & F1--F10 \\
Alt/Opt destro + 1--9 & Seleziona lo slot 1--9 \\
Alt/Opt destro + Left / Right & Slot precedente / successivo \\
Alt/Opt destro + Tab & Escape \\
Alt/Opt destro + Up / Down & PageUp / PageDown \\
Alt/Opt destro + Backspace & Conferma la chiusura di tutti i terminali \\
Alt/Opt destro + Space & Settings \\
Alt sinistro + Space & Tastiera Unicode \\
Alt/Opt destro + C / V & Copia selezione / incolla \\
\end{keytable}

Ctrl+Shift+B viene inoltrato senza modifiche ai programmi, incluso Emacs.

Sul \key{Type Folio statunitense}, Alt/Opt destro+\key{0} inserisce \code{+}.
Alt/Opt destro con il \key{tasto meno subito a sinistra di Backspace} inserisce
\code{=}. Con il metodo di input disattivato, Shift+6 produce un circonflesso
letterale: \code{c^2} rimane \code{c^2}. Non ci sono scorciatoie per lo zoom.

Le tastiere USB mantengono i normali tasti funzione. La scorciatoia Opt separata
usa Meta quando la tastiera USB lo fornisce.

Sono disponibili al massimo \key{nove slot indipendenti}. Selezionare uno slot
vuoto apre una shell. L'indicatore conta i terminali aperti: se sono aperti solo
1 e 9, il 9 mostra \key{Terminal 2 / 2 · slot 9}. Chiudere una shell libera
lo slot senza rinumerare le scorciatoie.

Ogni terminale ha shell, directory, composizione di input, immagini e 500 righe
di cronologia proprie. Le impostazioni di visualizzazione e input sono comuni.

\topic{Impostazioni e schermo}
Tocca \key{Settings}, a destra della barra, oppure premi
\key{Alt/Opt destro+Space}; funziona anche con la barra nascosta.
\key{Alt sinistro+Space} apre la tastiera Unicode.
Una scelta piena indica lo stato attivo; il bordo tratteggiato indica il focus.
Tab sposta il focus, frecce e Invio cambiano la scelta.

\key{Caps Lock:} scegli Control (predefinito) o Caps Lock.
\key{Dimensione:} pizzica con due dita oppure usa meno/più nelle impostazioni.
Movimenti lenti regolano un pixel alla volta, da \key{6 a 48 pixel}.

\key{Text darkness:} 50\% corrisponde alla resa originale. Valori più alti
scuriscono i bordi dei caratteri. \key{Minimum contrast} separa la luminanza
di testo e sfondo quando un'app sceglie colori simili; il valore iniziale è 35\%.

\begin{choices}
\key{Modalità e-paper} & \key{Compromesso} \\ \midrule
Fast & Attesa minima; forma d'onda per animazioni; più immagini residue \\
Balanced & Forma d'onda UI; raggruppamento dell'output in 32 ms \\
Crisp (predefinita) & Grigi più puliti; forma d'onda per contenuti; più lenta \\
Mono & Bianco e nero veloce; elimina i grigi dalle immagini \\
Saver & Raggruppamento in 120 ms; meno aggiornamenti durante le raffiche \\
\end{choices}

Si ridisegnano solo le righe cambiate, ma il pannello richiede comunque tempo.
Il risparmio di batteria di Saver non è stato misurato. Gli aggiornamenti
conservano una modalità salvata esplicitamente.

\begin{notice}
Tutte le impostazioni rimangono in \key{RAM} e si salvano insieme all'uscita
normale da Inkline. Sollevare il dito, chiudere Settings o sospendere non
scrive sulla flash. Una sessione senza cambiamenti non scrive nulla.
Un arresto anomalo o forzato, o la perdita di alimentazione, può perdere le modifiche non salvate.
\end{notice}

La barra inferiore legge la batteria una volta al minuto dall'interfaccia di
alimentazione Linux in sola lettura. \code{inkline-battery} mostra percentuale
e stato di carica; \code{inkline-battery --percentage} o \code{~/inkline battery}
mostra il valore breve usato nel prompt Bash senza colori di Inkline.

\topic{Unicode e metodi di input}
\key{Alt sinistro+Space} apre la tastiera Unicode. Tocca un carattere per inserirlo;
la tastiera rimane aperta. Tocca \key{Done} o premi Escape per chiuderla.

Le categorie includono lettere, segni combinanti, numeri, punteggiatura,
matematica/valute, simboli, spazi, caratteri di formato e uso privato.
Trascina verticalmente oppure usa frecce, PageUp/PageDown, Home/End e rotella.
\key{Tab} cambia categoria; Invio inserisce il carattere selezionato.

Scrivi un valore Unicode esadecimale, per esempio \code{03bb}, poi Invio per
inserire λ. Backspace modifica il valore. Il cerchio tratteggiato dei segni
combinanti è solo un aiuto visivo; si inserisce soltanto il segno. Spazi e
caratteri di formato hanno etichette. Il catalogo usa le tabelle Unicode di Qt,
escludendo controlli, valori non assegnati, surrogati e non caratteri. Alcuni glifi
richiedono font assenti; i valori privati richiedono un font adatto.

Nel selettore, premi \key{F2} o tocca \key{Settings}, poi \key{Input method}:
\begin{choices}
Off — direct keyboard & Digitazione statunitense normale \\
Romaji & Hiragana giapponese da input romanizzato \\
Pinyin & Cinese tramite pinyin \\
Zhuyin & Cinese con la mappa zhuyin di base \\
Wubi 86 & Cinese tramite codici di forma \\
US-International & Accenti latini comuni \\
\end{choices}

Scegli un risultato nella barra dei candidati. Seleziona
\key{Off — direct keyboard} nello stesso menu per tornare all'input statunitense.
La scelta vale per tutti i terminali e viene salvata all'uscita da Inkline.

\topic{Tocco, penna e appunti}
\key{Due dita su/giù:} scorrono la cronologia. Trascina verso il basso per
l'output precedente e verso l'alto per tornare al prompt. Ogni terminale
conserva 500 righe fisiche in RAM, oltre alla schermata corrente.

\key{Due dita a sinistra/destra:} cambiano uno slot per gesto.
Solleva le dita prima di cambiare ancora. Allargale o avvicinale per
modificare la dimensione del testo.

\key{Un dito:} trascina per selezionare. Alt/Opt destro+C copia;
Alt/Opt destro+V incolla. Tutti i terminali condividono appunti in RAM
fino a 128 KiB.

\key{Penna:} le app che attivano la segnalazione del mouse ricevono pressione,
movimento e rilascio come eventi del pulsante sinistro, con il protocollo
negoziato, incluso SGR. Negli altri casi il trascinamento seleziona testo.
Tieni premuto \key{Shift} per forzare la selezione locale nelle app che usano il mouse.

\key{OSC 52} permette ai programmi del terminale, anche tramite SSH, di
leggere o sostituire gli appunti di Inkline, separati da quelli dei quaderni.
Alt/Opt destro+X copia e spiega che l'eliminazione deve avvenire nell'editor:
OSC 52 non può cancellare il testo di un editor.

\key{Grafica:} i trasferimenti kitty inline e tramite memoria condivisa restano
in RAM. Le modalità file e file temporaneo sono disabilitate. Sixel funziona
su richiesta esplicita, ma non viene annunciato automaticamente. I segnaposto
kitty, le animazioni e alcune modalità sixel sono ancora incompleti.

\code{clear} e la cancellazione dell'intero schermo rimuovono testo e immagini
visibili. Le immagini fuori schermo vengono liberate quando la memoria scarseggia.
Non esiste una cache su disco; le immagini grandi consumano comunque la RAM limitata.

\topic{Tastiera USB e macchina da scrivere}
Apri \key{Settings → USB (pagina 3)}. \key{Network} ripristina Ethernet USB;
\key{Send keyboard} trasforma il tablet in tastiera per un computer;
\key{Receive keyboard} riceve una tastiera tramite adattatore OTG alimentato.
Type Folio resta disponibile. Installa l'utilità Python del sito per inviare
pressioni di tasti; non servono moduli Python aggiuntivi.

Scegli Send keyboard, collega USB-C, scegli il profilo e apri
\key{Open typewriter}. Se Send keyboard è attivo, l'invio dal vivo parte
automaticamente. In pausa, l'editor mostra chiaramente
\key{PAUSED · USB OUTPUT OFF}; l'invio attivo ha un indicatore distinto.
\key{Escape} o \key{Stop}
ferma l'invio e rilascia i tasti. Anche uscire dalla schermata, cambiare terminale,
perdere il focus, aprire un'app nativa o sospendere mette in pausa.
Dopo la riconnessione non viene reinviato nulla automaticamente.

\begin{choices}
US / ASCII & Layout US, Bloc Maiusc disattivato. Rifiuta testo non ASCII. \\
Mac & Sorgente Unicode Hex Input, Bloc Maiusc disattivato. Fino a U+FFFF;
rifiuta caratteri dei piani supplementari. \\
Linux & Layout US, Bloc Maiusc disattivato. L'app deve accettare Ctrl+Shift+U,
codice esadecimale, Invio. Non tutte lo supportano. \\
Windows & Installa WinCompose sul PC, imposta Compose su Right Alt e attiva
l'immissione Unicode. Layout US, Bloc Maiusc disattivato. Invia Compose, u,
sei cifre esadecimali, Invio. \\
\end{choices}
WinCompose: \url{https://wincompose.info/}. Le regole personali non devono sostituire
questa sequenza. HID non offre un canale universale di testo Unicode; contano
configurazione e font del destinatario. Prova prima un breve testo.

\key{Input method} offre Off/US diretto, Romaji, Pinyin, Zhuyin, Wubi 86 e
US-International. Invia solo i caratteri confermati. Cambiare metodo mette in
pausa: riprendi con Start. \key{Off} ripristina la normale tastiera US.
\key{Unicode} o Alt sinistro+Space apre il selettore di caratteri per l'uscita USB.

Typewriter è un editor locale: il documento resta in RAM mentre Inkline è aperto.
La digitazione modifica il documento e, con la modalità live attiva, invia ogni
carattere confermato al computer. Copia, taglia, incolla, selezione e scorrimento
agiscono nell'editor locale. \key{Files \& send} carica e salva UTF-8 con qualsiasi
nome, sceglie 5--80 caratteri al secondo, invia l'intero documento e interrompe
l'invio. Solo Save o l'uscita normale scrivono il documento. Un documento senza
nome usa \path{~/.inkline-typewriter-draft} all'uscita per il recupero.
\key{Exit typewriter} torna alle impostazioni USB.

\topic{Inviare file e ripristinare la rete}
I comandi \code{inkline-usb} e \code{keyboard-send} sono in
\path{~/.local/bin}. \code{~/inkline usb}, \code{~/inkline type} e
\code{inkline-type} restano alias compatibili. Gli strumenti Outpost rimangono intatti.
\UsbExamples
I file UTF-8 \key{non hanno un limite fisso di dimensione}. Si leggono piccoli
blocchi senza conservare il documento intero in RAM o creare file temporanei.
I file riposizionabili vengono verificati interamente prima dell'invio e poi
riletti: non modificare la sorgente durante l'operazione. Le pipe si verificano
man mano; un errore può seguire testo già inviato. Il BOM UTF-8 iniziale viene
rimosso e CRLF/CR diventano nuove righe. Invio e Tab agiscono sull'app attiva.

\code{--cps} (predefinito 20) o \code{--delay-ms} regola la velocità; Unicode
richiede più rapporti. \code{--start-delay} lascia tempo per spostare il focus.
Dry-run mostra i rapporti senza USB. Ctrl+C o \code{--stop} annulla e rilascia
i tasti. Funziona un solo mittente alla volta. Una coda interattiva piena mette
in pausa. Controlla il documento prima di riprovare un invio interrotto.
Testo e stato temporaneo non scrivono sulla memoria flash. Il profilo si salva
con le altre impostazioni all'uscita normale.

\key{Network} o \code{inkline-usb network} ripristina Ethernet. Anche l'uscita da
Inkline, un arresto anomalo o un riavvio di emergenza ripristinano la rete.
La modalità USB non viene salvata come preferenza di avvio.
\UsbRecoveryExample
Il timer è indipendente dalla shell. Cambia modalità sul tablet o tramite SSH
Wi-Fi, mai tramite SSH USB. In caso di errore si tenta il ripristino di Ethernet.
Una sessione PPP satellitare attiva o un gadget sconosciuto mantiene la porta;
vale anche per la pulizia all'uscita. Opt+RightAlt+Backspace riavvia Inkline,
ma chiude tutti i terminali e perde il lavoro non salvato.

\topic{Avviare programmi con scorciatoie}
Registra le scorciatoie da una shell Inkline. Gli argomenti sono letterali;
chiama esplicitamente una shell per usarne la sintassi. I programmi avviati
leggono \code{~/.bashrc} tramite Bash.

Settings → Command shortcuts (pagina 2) permette di assegnare un comando a una lettera. Scegli Terminal o Native e-paper app, poi Save. Remove chiede conferma; Reset draft scarta le modifiche. T è bloccata. Solo Save e la rimozione confermata scrivono sullo spazio di archiviazione. Lo strumento da riga di comando gestisce le stesse associazioni.

Tutte le scorciatoie globali usano ora Opt+RightAlt, incluse T e il recupero tenendo premuto Backspace. Usa il tasto Opt separato del Folio; sulle tastiere USB usa Windows/Command (Super). Ctrl+Alt resta disponibile per Emacs. Chiudi e riapri Inkline dopo l’aggiornamento; l’installazione conserva i terminali aperti.

\begin{shell}
~/inkline shortcut register e emacs
~/inkline shortcut register p python3
~/inkline shortcut list
~/inkline shortcut deregister p
\end{shell}

\key{Opt+RightAlt+E} apre Emacs in uno slot libero.
\code{~/inkline run PROGRAM ARG...} avvia senza registrare una scorciatoia.
Una nuova registrazione sostituisce l'associazione; rimuoverla non tocca i
programmi in esecuzione. I tasti seguono le posizioni fisiche statunitensi.
Metti la punteggiatura fra virgolette e usa percorsi completi per programmi
personalizzati, che devono essere già installati.

Per un'app Qt/e-paper nativa che disegna il proprio schermo:
\begin{shell}
~/inkline shortcut register a --epaper /home/root/my-app
\end{shell}
Un'app nativa alla volta mette in pausa Inkline e le sue shell. Il lanciatore
fornisce impostazioni Qt e cache in RAM, scarta i log e riprende Inkline alla
chiusura. Le app devono usare l'input Qt senza catturare la tastiera in esclusiva.

\begin{notice}
\key{Opt+RightAlt+T è permanente.} Non si può registrare né rimuovere.
Ferma l'app nativa e torna ai terminali esistenti. Non tenere premuto Ctrl per questa scorciatoia;
la rimappatura di Caps Lock vale soltanto dentro Inkline.
\end{notice}

\key{Emergenza:} tieni Opt+RightAlt+Backspace per due secondi per fermare l'app
e i suoi discendenti e riavviare Inkline. Tutti i terminali si chiudono e il
lavoro non salvato si perde. Un'app root può impedire il recupero catturando
l'input, cambiando servizi o esaurendo risorse. Il riavvio del dispositivo è
l'ultima possibilità.

\topic{Installazione, aggiornamenti e guida}
Il pacchetto è per \key{reMarkable 2 con firmware 3.27} e usa le librerie
Qt 6.8 e-paper di sistema. Download e verifica delle somme di controllo:

\InstallLink

L'installer controlla hardware, firmware, font, avvio di Qt, directory
temporanee in RAM e swap disabilitato. Le versioni sono sotto
\path{/home/root/.local/share/inkline}; il lanciatore da tastiera parte all'avvio.
Gli aggiornamenti conservano i terminali aperti. \key{Chiudi e riapri Inkline
per usare la nuova versione.} Le versioni precedenti restano per tornare indietro
e occupano spazio aggiuntivo.

Questo unico PDF comprende tutte e nove le lingue ed è incluso nel pacchetto.
L'installazione tenta l'importazione in \key{My files} tramite l'interfaccia
web USB, quando è abilitata e i quaderni sono aperti. Un'importazione riuscita
viene registrata per non duplicare lo stesso manuale.

Se l'importazione automatica non è disponibile, collega USB, chiudi Inkline
e abilita \key{USB web interface} nelle impostazioni di archiviazione dei quaderni.
Da SSH esegui:
\begin{shell}
~/inkline manual
\end{shell}
Puoi anche scaricare il PDF e importarlo dal browser USB o dall'app reMarkable
desktop/mobile. L'importatore non modifica mai direttamente il database dei
quaderni né interrompe i terminali. Il manuale importato resta in My files
dopo la disinstallazione di Inkline.

Da SSH, torna ai quaderni o disinstalla Inkline con:
\begin{shell}
~/inkline stop
~/inkline uninstall
\end{shell}
Le istruzioni per il software facoltativo, inclusa l'installazione LaTeX
consigliata, rimangono sul sito.

\topic{Funzionamento, memoria e spazio}
Qt invia i tasti a un mappatore comune per scorciatoie e Caps Lock.
libghostty-vt codifica i tasti del terminale; ogni slot collega il programma
a uno pseudoterminale (PTY) separato. Le shell ricevono
\code{HOME=/home/root}; Bash interattivo legge \code{~/.bashrc}.

L'output passa nell'analizzatore del terminale e nel decoder sixel. libghostty
gestisce celle, cursore, schermate alternative, cronologia e immagini kitty.
Il renderer conserva un'immagine grigia e ridisegna le righe cambiate; il backend
e-paper Qt aggiorna il pannello. Un demone evdev indipendente gestisce scorciatoie
globali e accensione; systemd gestisce sospensione e ripresa.

Ogni terminale limita le allocazioni del nucleo a 64 MiB, le immagini kitty
a 16 MiB per schermata e le immagini sixel conservate a 16 MiB. Rendering e
app richiedono memoria aggiuntiva. Gli slot vuoti non allocano terminali.
Nove app grandi possono superare la RAM disponibile.

Immagini temporanee, cache e sessioni usano RAM. Software installato, documenti,
impostazioni salvate all'uscita e file scritti dalle app usano memoria persistente.
Inkline non impedisce ai programmi di scrivere file.

Il pacchetto Python facoltativo usa Python standard. Per ridurre le scritture
di cache, aggiungi a \code{~/.bashrc} sul tablet:
\begin{shell}
export PYTHONDONTWRITEBYTECODE=1
export PIP_NO_CACHE_DIR=1
\end{shell}
Esegui \code{source ~/.bashrc} o apri una nuova shell. Usa
\code{python3 -m pip install --no-compile PACKAGE} per evitare la compilazione
bytecode aggiuntiva effettuata da pip durante l'installazione. La modalità
isolata di Python ignora le sue variabili d'ambiente. La guida completa è sul sito.

Inkline è software GPL-3.0-or-later. I pacchetti includono sorgenti corrispondenti
e avvisi sulle dipendenze. I resoconti distinguono test automatici e comportamenti
fisici ancora da confermare sul dispositivo.
